\section{การบำรุงรักษาและการแก้ไขปัญหาเบื้องต้น}
\label{sec:troubleshooting}

เพื่อให้เครื่องสเปกโทรโฟโตมิเตอร์ทำงานได้อย่างมีเสถียรภาพและมีความแม่นยำสูงสม่ำเสมอ ผู้ใช้งานควรปฏิบัติตามแนวทางการบำรุงรักษาและการแก้ไขข้อขัดข้องเบื้องต้นดังต่อไปนี้

\subsection{การแก้ไขข้อขัดข้องที่พบบ่อย}
ปัญหาที่อาจเกิดขึ้นระหว่างการประกอบ การติดตั้ง หรือการใช้งาน และแนวทางการแก้ไขได้สรุปไว้ในตารางที่ \ref{tab:troubleshooting}

\begin{table}[H]
    \caption{ปัญหาและแนวทางการแก้ไขข้อขัดข้องเบื้องต้น}
    \label{tab:troubleshooting}
    \begin{tabularx}{\textwidth}{p{3.8cm}p{4.2cm}X}
        \toprule
        \textbf{อาการที่พบ} & \textbf{สาเหตุที่เป็นไปได้} & \textbf{แนวทางการแก้ไข} \\
        \midrule
        หน้าจอเป็นสีขาวสว่างค้าง & มีการเรียกใช้ไลบรารี TFT\_eSPI ซ้ำซ้อน & ลบหรือย้ายโฟลเดอร์ TFT\_eSPI ทั่วไปออกจาก Arduino/libraries และใช้ Seeed\_Arduino\_LCD เท่านั้น \\
        \midrule
        กดปุ่มแล้วหลอด LED ไม่สว่าง & สายสัญญาณหลวม หรือโปรแกรมขับขาผิด & ตรวจสอบการเสียบสาย Grove พอร์ตซ้าย (D0) และใช้คำสั่งขับสัญญาณแบบ NeoPixel \\
        \midrule
        เซนเซอร์ไม่ตรวจพบข้อมูล & สายสัญญาณ I2C หลุด หรือหน้าสัมผัสไม่แน่น & เสียบสายเซนเซอร์เข้าพอร์ต Grove ด้านขวาให้แน่นหนา และตรวจดูที่อยู่ I2C (0x29) \\
        \midrule
        ค่าการวัดผันผวนผิดปกติ & หลอดคิวเวตต์มีรอยนิ้วมือหรือมีแสงภายนอกรั่ว & เช็ดทำความสะอาดผนังหลอดคิวเวตต์ด้วยกระดาษเช็ดเลนส์ และปิดฝาครอบกันแสงให้สนิท \\
        \midrule
        ไมโครคอนโทรลเลอร์ไม่เชื่อมต่อคอมพิวเตอร์ & พอร์ตสื่อสาร USB ค้าง & โยกสวิตช์เปิดปิดด้านข้างลง 2 ครั้งติดกันเพื่อเข้าโหมด Bootloader แล้วอัปโหลดใหม่อีกครั้ง \\
        \bottomrule
    \end{tabularx}
\end{table}

\subsection{ข้อแนะนำการบำรุงรักษาระบบเชิงทัศนศาสตร์}
\begin{tipbox}
    \begin{enumerate}
        \item หลีกเลี่ยงการใช้นิ้วมือสัมผัสบริเวณผนังด้านโปร่งใสของหลอดคิวเวตต์ ให้จับเฉพาะส่วนขอบด้านบนหรือด้านที่มีผิวฝ้าเท่านั้น
        \item ก่อนใส่หลอดคิวเวตต์ลงในช่องวัด ควรตรวจดูว่าไม่มีหยดน้ำหรือคราบสารละลายเกาะอยู่ที่ผิวด้านนอก เพื่อป้องกันการกัดกร่อนเซนเซอร์และหลอดไฟ
        \item หลังเสร็จสิ้นการตรวจวัดในแต่ละวัน ควรนำหลอดคิวเวตต์ออกมาล้างทำความสะอาดด้วยน้ำกลั่นและคว่ำให้แห้งสนิทก่อนจัดเก็บ
        \item ควรเก็บรักษาตัวเครื่องไว้ในกล่องกันกระแทกที่แห้งและพ้นจากแสงแดดจัด เพื่อรักษาเสถียรภาพของชิ้นส่วนที่พิมพ์ด้วยเครื่องพิมพ์สามมิติ
    \end{enumerate}
\end{tipbox}
